\documentclass[12pt]{article}
\usepackage{geometry,fancyhdr,xr,hyperref,ifpdf,amsmath,rcs,shortcuts}
\usepackage{lastpage,longtable,color,paunits,amsmath}
\geometry{letterpaper,headsep=5mm,head=10mm,hmargin={20mm,30mm},bottom=10mm,top=5mm}


\pagestyle{fancy} 

\RCS $Revision: 1.5 $
\RCS $Date: 2002/01/09 03:50:54 $

\fancypagestyle{first}{
\lhead{}
\chead{Wet bulb vs $\theta_e$}
\rhead{page~\thepage/\pageref{LastPage}}
\lfoot{} 
\cfoot{} 
\rfoot{}
}

\ifpdf
    \usepackage[pdftex]{graphicx} 
    \usepackage{hyperref}
    \pdfcompresslevel=0
    \DeclareGraphicsExtensions{.pdf,.jpg,.mps,.png}
\else
    \usepackage{hyperref}
    \usepackage[dvips]{graphicx}
    \DeclareGraphicsRule{.eps.gz}{eps}{.eps.bb}{`gzip -d #1}
    \DeclareGraphicsExtensions{.eps,.eps.gz}
\fi
\externaldocument[lec6,]{lec6}

\begin{document}

\pagestyle{first}

\begin{center}
Atsc405: Lecture Summary: Monday, Jan. 30\\
\end{center}




\section{Converting between $\theta_e$ and $\theta_w$}

Note that Canadian and  UK  tephigrams  do not use $\theta_e$ to label their moist
adiabats, they use $\theta_w$, the web bulb potential temperature, defined
as the temperature found by saturated adiabatic ascent/descent to 100 kPa
(i.e.~the name comes from the fact that this is what you would get
if you physically moved a saturated wick upward/downward in an adiabatic
parcel.

Convince yourself that from the above definition we get the following limits
of integration for a saturated parcel with a temperature of $\theta_w$ 
and mixing ratio $w_{s0}$:

\begin{equation}
\label{eq:thetaodew}
\int_{\theta_w}^{\theta_e} d \ln \theta = -\int_{w_{s0}}^0 \frac{l_v}{c_p T} d w_s \approx - 
\int_{(w_{s0},\theta_w)}^0 d \left ( \frac{l_v w_s}{c_p T} \right )  
\end{equation}
which integrated gives;

\begin{equation}
  \label{eq:thetaint}
  \theta_e = \theta_w \exp \left ( \frac{l_v w_{s0}}{c_p \theta_w} \right )
\end{equation}

if $ \frac{l_v w_{s0}}{c_p \theta_w} \ll 1$ (is it?) then we can expand

\begin{equation}
  \label{eq:exp}
  \exp(x) = 1 + x + \frac{x^2}{2} + \ldots
\end{equation}
and drop order 2 and higher to get:

\begin{equation}
  \label{eq:expand}
  \theta_e=\theta_w \left ( 1 + \frac{l_v w_{s0}}{c_p \theta_w} \right )
\end{equation}

or rearranging:

\begin{equation}
  \label{eq:expand2}
  \theta_e - \theta_w  \approx  \frac{l_v w_{s0}}{c_p}
\end{equation}

Try this out for the point ($p=100\ \un{kPa},T=10\degc = 283.15\ K, w_{s}= 7.8\ \un{g\,kg^{-1}}$),
use lecture 6 $\theta_e$ definition (\ref{lec6,eq:thetae2}) and find for the $10\degc$ $\theta_w$
adiabat that
$\theta_e= 283.15 \exp \left ( \frac{2.5 \times 10^6 \times 0.0078}{ 1004 \times 283.15} \right )$
$=303.25K$
so that $\theta_e - \theta_w=303.25 K - 283.15 K =20.1$ degrees.  Compare this to the
prediction of (\ref{eq:expand2}) that 
$\theta_e - \theta_w = \frac{2.5 \times 10^6 \times 0.0078}{ 1004}$=19.4 degrees.

Note the technical difficulty if we want to do any better than this (if, for example,
we need more exact thermodynamics for a computer model) -- getting $\theta_w$ in
terms of $\theta$, $T_{lcl}$, and $w_s$ via (\ref{eq:thetaint}) is a bit of a mess.

\end{document}


%%% Local Variables:
%%% mode: latex
%%% TeX-master: t
%%% End:
